\documentclass{article}
\usepackage{amsmath, amssymb, amsthm}
\usepackage{hyperref}
\usepackage{geometry}
\geometry{margin=1in}

\title{Erd\H{o}s Problem \#596}
\date{Last updated: 2025-08-31}
\author{Source: erdosproblems.com}

\begin{document}
\maketitle

\noindent\textbf{Status:} OPEN \quad \textbf{No prize}

\noindent\textbf{Tags:} graph theory, ramsey theory, set theory

\medskip
\noindent\textbf{Problem Statement:}

For which graphs $G_1,G_2$ is it true that
{UL}
{LI} for every $n\geq 1$ there is a graph $H$ without a $G_1$ but if the edges of $H$ are $n$-coloured then there is a monochromatic copy of $G_2$, and yet{/LI}
{LI} for every graph $H$ without a $G_1$ there is an $\aleph_0$-colouring of the edges of $H$ without a monochromatic $G_2$.
{/UL}

\bigskip
\noindent\textbf{Remarks:}

Erd\H{o}s and Hajnal originally conjectured that there are no such $G_1,G_2$, but in fact $G_1=C_4$ and $G_2=C_6$ is an example. Indeed, for this pair Ne\v{s}et\v{r}il and R\"{o}dl established the first property and Erd\H{o}s and Hajnal the second (in fact every $C_4$-free graph is a countable union of trees). 

Whether this is true for $G_1=K_4$ and $G_2=K_3$ is the content of [595].

\bigskip
\noindent\small{Source: \url{https://www.erdosproblems.com/596}}
\end{document}
